\documentclass[11pt]{article}
\usepackage[margin=0.82in]{geometry}
\usepackage{amsmath,amssymb,mathrsfs,lmodern}
\usepackage[T1]{fontenc}
\usepackage{needspace}
\usepackage{longtable,booktabs,array,calc}
\usepackage[colorlinks=true,linkcolor=blue,urlcolor=blue]{hyperref}
\setlength{\emergencystretch}{3em}
\setlength{\parskip}{0.3em}
\setcounter{secnumdepth}{0}
\providecommand{\tightlist}{\setlength{\itemsep}{0pt}\setlength{\parskip}{0pt}}
\title{Complete splitting of planar Robin eigenvalues by a single shrinking hole}
\author{Henry Zweiman}
\date{September 16, 2026}
\begin{document}
\maketitle
\textbf{AI-assisted research preprint, version 1.0.} Prepared with
OpenAI Codex. The proofs and literature comparison were checked by the
originating assistant; they have not been independently verified by
human experts or peer reviewed. Source-access limits are stated in
Section 20.

\Needspace{8\baselineskip}\section{Abstract}\label{abstract}

We study planar Robin eigenvalues when a scaled interior hole shrinks to
a prescribed point, with the same fixed nonzero real Robin parameter on
the outer and inner boundaries. We establish a complete multiscale
description of every branch of a multiple limiting eigenvalue for
arbitrary Lipschitz holes with connected exterior. The leading
coefficients are Schur complements of an exterior harmonic energy matrix
on the vanishing-order filtration. A nonnodal direction, when present,
moves at first order with the sign of the Robin parameter; all nodal
directions move down at even powers determined by their orders. We then
construct an area-preserving analytic family of strictly convex holes
near a disk. Outside a countable exceptional set of shape parameters, a
single selected shape separates every fixed eigenvalue cluster,
including at common nodal points. The smallness threshold for the hole
size depends on the spectral window. Quantitative correctors, a two-term
nodal formula and a nonnodal logarithmic correction accompany the
splitting result.

\Needspace{8\baselineskip}\section{1. Introduction and
attribution}\label{introduction-and-attribution}

Small holes constitute a singular domain perturbation. A complete
splitting theorem needs estimates that remain effective for
eigenfunctions vanishing to different orders at the concentration point;
one leading unweighted matrix estimate can miss entire branches. In the
plane, the exterior Neumann corrector also requires a normalization
modulo constants because a homogeneous Sobolev estimate does not control
them.

The main result is stated in Section 2. Its proof uses the
arbitrary-hole planar asymptotics in Sections 3-15 and an explicit
variation of the exterior energy coefficients in Sections 16-19. The
hole center remains fixed. The selected hole shape is independent of the
spectral window, whereas its size threshold may depend on that window.
We do not claim a uniform threshold for the entire infinite spectrum.

Felli-Roychowdhury-Siclari {[}FRS26{]} establish Robin spectral
stability for either sign of the parameter and sharp simple-eigenvalue
asymptotics, with a planar restriction to round holes for their nodal
coefficient. Their final Remark 2.10 leaves multiple Robin eigenvalues
and splitting for future work. Their perturbative framework, uniform
extension and trace estimates are prior work. We include the needed
compactness arguments and derive quantitative cutoff residuals to make
the planar remainder explicit.

Felli-Liverani-Ognibene {[}FLO26{]} obtain generic reduction of
multiplicity under removal of a spherical Neumann hole as its center
varies, and complete splitting for double eigenvalues. Their
arbitrary-hole leading matrix is proved in dimensions at least three;
their planar theorem treats disks. Their earlier work {[}FLO25{]}
supplies the related simple-eigenvalue analysis.

Abatangelo-Léna-Musolino {[}ALM26{]} correct the multiple-eigenvalue
coefficients in the Dirichlet-hole theory of {[}ALM22{]}. Their
corrected vanishing-order and matrix analysis is essential prior work:
the projection method, filtration, and Schur-complement interaction are
not claimed new here. The present calculations establish the Robin PDE
coefficients and handle the possible coexistence of positive and
negative shifts.

Shape differentiation of generalized polarization tensors and its
Fourier interpretation are also prior work {[}AKLZ12{]}. Full splitting
for other singular perturbations occurs in
Gadyl\textquotesingle shin\textquotesingle s boundary-patch problem
{[}Gad93{]}. The distinction pursued here is the fixed-center,
single-hole Robin conclusion with arbitrary finite multiplicity and
explicit behavior at all nodal orders. Earlier expansion and
repeated-eigenvalue frameworks {[}AKLZ10, LdC12, Dab17, Dab22{]} remain
relevant. The inspected report {[}Dab17{]} gives leading matrix laws for
grounded and conductivity inclusions; its stated applications do not
give the fixed-center all-order Robin conclusion. The source-access
limits of this comparison are disclosed in Section 20.

\Needspace{8\baselineskip}\section{2. Main splitting theorem}\label{main-splitting-theorem}

Fix a bounded connected Lipschitz planar domain \(\Omega\), an interior
point \(x_0\), and a nonzero real Robin parameter \(\alpha\). Translate
\(x_0\) to zero. Put

\[
 f(\theta)=\sum_{n=1}^{\infty}e^{-n^2}\cos(2n\theta),
 \qquad
 \widetilde\Sigma_t=\{(r,\theta):0\le r<1+t f(\theta)\},
 \qquad
 \Sigma_t=\left(\frac{\pi}{|\widetilde\Sigma_t|}\right)^{1/2}
                      \widetilde\Sigma_t.
 \tag{H1}
\]

For all \(|t|<t_*\), with an absolute sufficiently small \(t_*>0\),
these are strictly convex real-analytic, centrally symmetric holes of
area \(\pi\). They contain zero and have connected exteriors. They
converge to the unit disk in every fixed \(C^q\) norm as \(t\to0\).
Indeed, the Fourier series and all its derivatives converge on each
complex strip, \(f\) has zero average, and strict positivity of
curvature persists in a small \(C^2\) neighborhood of the circle. The
area normalization factor is \(1+O(t^2)\).

\Needspace{6\baselineskip}\textbf{Theorem 2.1 (complete fixed-center splitting).} There is a
countable set \(\mathcal N\subset(-t_*,t_*)\), depending on
\(\Omega,x_0,\alpha\), such that the following holds. For every
\(t\notin\mathcal N\) and every eigenvalue \(\lambda\) of the
unperforated Robin problem, there is \(\varepsilon_\lambda(t)>0\) such
that every branch converging to \(\lambda\) is simple for

\[
 0<\varepsilon<\varepsilon_\lambda(t)
\]

in \(\Omega\setminus\varepsilon\overline\Sigma_t\), with the same Robin
parameter on both boundaries. In particular, for each finite \(M\),
there is \(\varepsilon_M(t)>0\) such that the first \(M\) perturbed
eigenvalues are simple whenever \(0<\varepsilon<\varepsilon_M(t)\).

The center is fixed and may be a common nodal point of an eigenspace.
The shape is fixed before \(\varepsilon\to0\), and can be chosen
arbitrarily close to a disk. The threshold can depend on the spectral
window. \textbf{No common positive threshold for the entire infinite
spectrum is asserted.} The theorem is not uniform in \(\alpha\), in the
center, or in the outer domain. For a single eigenvalue cluster, every
sufficiently small nonzero \(t\) works, rather than merely all \(t\)
outside a countable set.

The finite-window assertion also holds with any prescribed area in place
of \(\pi\), by a fixed scaling. The unperforated simple eigenvalues
require no splitting argument: convergence and a spectral gap already
give an isolated one-dimensional perturbed cluster.

\Needspace{8\baselineskip}\section{3. Geometry and the simple nodal
branch}\label{geometry-and-the-simple-nodal-branch}

Let \(\Omega\subset\mathbb R^2\) be a bounded connected Lipschitz domain
containing zero. Let \(\Sigma\) be a nonempty bounded Lipschitz open set
with \(\overline\Sigma\subset B_R\), and suppose
\(B_R\setminus\overline\Sigma\) is connected. We do not require that
zero lie in \(\Sigma\). Put

\[
 D=\mathbb R^2\setminus\overline\Sigma,
 \qquad \Omega_\varepsilon=\Omega\setminus\varepsilon\overline\Sigma.
\]

Fix \(\alpha\in\mathbb R\). Impose \(\partial_\nu u+\alpha u=0\) on
every boundary component. The normal on a hole points into the hole. Let
\(\lambda\) be a simple eigenvalue in \(\Omega\), with real
eigenfunction \(u\) normalized by \(\|u\|_{L^2(\Omega)}=1\). Assume
\(u(0)=0\), and write its interior Taylor expansion as

\[
 u(x)=P_k(x)+O(|x|^{k+1}),\qquad
 \nabla u(x)=\nabla P_k(x)+O(|x|^k),\qquad k\ge1.
\]

Here \(P_k\) is a nonzero homogeneous harmonic polynomial. Let \(W_k\)
be the exterior harmonic function determined by

\[
 \partial_\nu W_k=\partial_\nu P_k\quad\hbox{on }\partial\Sigma,
 \qquad W_k(y)=O(|y|^{-1})\quad (|y|\longrightarrow\infty).
 \tag{S1}
\]

Define

\[
 E_k=\int_D|\nabla W_k|^2,
 \qquad C_k=E_k+\int_\Sigma|\nabla P_k|^2.
 \tag{S2}
\]

\Needspace{6\baselineskip}\textbf{Theorem 3.1 (simple nodal branch).} Let \(\lambda_\varepsilon\)
be the eigenvalue converging to \(\lambda\). Then

\[
 \lambda_\varepsilon
 =\lambda-C_k\varepsilon^{2k}
   +O(\varepsilon^{2k+1}).
 \tag{S3}
\]

In particular \(C_k>0\) and the eigenvalue decreases for sufficiently
small holes, for either sign of \(\alpha\). If \(u_\varepsilon\) is
normalized with positive overlap with \(u\), and
\(\chi\in C_c^\infty(\Omega)\) equals one near zero, then for every
\(\eta\in(0,1/2)\),

\[
 \left\|u_\varepsilon-u+
   \varepsilon^k\chi W_k(\cdot/\varepsilon)\right\|_{H^1(\Omega_\varepsilon)}
 =O_\eta(\varepsilon^{k+1-\eta}).
 \tag{S4}
\]

Consequently, for each fixed \(S\) with \(\overline\Sigma\subset B_S\),

\[
 \varepsilon^{-k}u_\varepsilon(\varepsilon\,\cdot)
 \longrightarrow P_k-W_k
 \quad\hbox{in }H^1(B_S\setminus\overline\Sigma),
 \tag{S5}
\]

and

\[
 \varepsilon^{-2k}\|u_\varepsilon-u\|_{H^1(\Omega_\varepsilon)}^2
 \longrightarrow E_k.
 \tag{S6}
\]

Constants can depend on the fixed geometry, eigenvalue, \(\alpha\),
cutoff and \(\eta\). No uniformity over families of domains or
eigenvalues is claimed.

\Needspace{8\baselineskip}\section{4. Uniform estimates and spectral
convergence}\label{uniform-estimates-and-spectral-convergence}

The facts in this section are established in {[}FRS26{]}. We include
their elementary scaling and compactness proofs to specify exactly which
geometry and parameter assumptions are needed. The fixed-domain trace,
extension, Sobolev and compact embedding theorems for bounded Lipschitz
domains are used in their usual forms.

Let \(A=B_R\setminus\overline\Sigma\) be connected, and fix a bounded
extension from \(H^1(A)\) to \(H^1(B_R)\) that agrees with the original
function on \(A\). For \(v\in H^1(\Omega_\varepsilon)\) apply this
extension after rescaling to \(v(\varepsilon\cdot)\) and subtracting its
mean on \(A\). Restore the mean after extension.
Poincaré\textquotesingle s inequality on \(A\) gives an extension
\(\mathscr E_\varepsilon v\in H^1(\Omega)\), equal to \(v\) off the
hole, such that

\[
 \int_{\varepsilon\Sigma}|\nabla\mathscr E_\varepsilon v|^2
       \le C\int_{\varepsilon A}|\nabla v|^2,
 \qquad
 \int_{\varepsilon\Sigma}|\mathscr E_\varepsilon v|^2
       \le C\int_{\varepsilon A}|v|^2
                  +C\varepsilon^2\int_{\varepsilon A}|\nabla v|^2.
 \tag{F1}
\]

The gluing across the hole boundary is valid because the traces agree.
Hence the extension is uniformly bounded in \(H^1\). The connectedness
hypothesis is used in the Poincaré estimate; no connectedness of
\(\Sigma\) itself is needed.

The fixed annular trace inequality, rescaled, yields

\[
 \|v\|_{L^2(\partial(\varepsilon\Sigma))}^2
 \le C\varepsilon^{-1}\|v\|_{L^2(\varepsilon A)}^2
        +C\varepsilon\|\nabla v\|_{L^2(\varepsilon A)}^2.
 \tag{F2}
\]

Given \(0<\eta<1/2\), apply planar Sobolev embedding to
\(\mathscr E_\varepsilon v\) with exponent \(p=2/\eta\) and then
Hölder\textquotesingle s inequality on \(\varepsilon A\). This gives

\[
 \|v\|_{L^2(\varepsilon A)}^2
    \le C_\eta\varepsilon^{2-2\eta}
                         \|v\|_{H^1(\Omega_\varepsilon)}^2,
 \qquad
 \|v\|_{L^2(\partial(\varepsilon\Sigma))}^2
    \le C_\eta\varepsilon^{1-2\eta}
                         \|v\|_{H^1(\Omega_\varepsilon)}^2.
 \tag{F3}
\]

A fixed outer-boundary collar misses the hole. The trace inequality on
this collar, with an arbitrarily small coefficient of the gradient term,
implies that for every \(\delta>0\)

\[
 \|v\|_{L^2(\partial\Omega)}^2
    \le\delta\|\nabla v\|_{L^2(\Omega_\varepsilon)}^2
       +C_\delta\|v\|_{L^2(\Omega_\varepsilon)}^2.
 \tag{F4}
\]

One can obtain this by multiplying by a fixed cutoff equal to one near
\(\partial\Omega\) and applying the compact trace inequality on
\(\Omega\); the cutoff vanishes near the hole. Combining (F3)-(F4) gives
constants \(c,a,b>0\), independent of sufficiently small
\(\varepsilon\), for which

\[
 a\|v\|_{H^1(\Omega_\varepsilon)}^2
 \le q_\varepsilon(v,v)
 :=\int_{\Omega_\varepsilon}(|\nabla v|^2+cv^2)
           +\alpha\int_{\partial\Omega_\varepsilon}v^2
 \le b\|v\|_{H^1(\Omega_\varepsilon)}^2.
 \tag{F5}
\]

This holds for either sign of the fixed real \(\alpha\). Choose the
constants so that the analogous form on \(\Omega\) is also coercive. The
associated operators are positive self-adjoint with compact resolvent.

\Needspace{6\baselineskip}\textbf{Lemma 4.1 (spectral convergence).} Each ordered Robin eigenvalue
of \(\Omega_\varepsilon\) converges, counting multiplicity, to the
corresponding eigenvalue of \(\Omega\).

\textbf{Proof.} For the upper bound in the minmax principle, restrict
the first \(j\) eigenfunctions of \(\Omega\) to \(\Omega_\varepsilon\).
They are analytic near zero, so the removed volume and inner-boundary
terms tend to zero uniformly on their finite-dimensional span. Their
mass matrix tends to the identity. Thus
\(\limsup\lambda_j(\Omega_\varepsilon)\le\lambda_j(\Omega)\).

For the reverse inequality take an orthonormal collection of the first
\(j\) perforated-domain eigenfunctions. Their shifted energies are
uniformly bounded by the upper bound and (F5). Their extensions are
uniformly bounded in \(H^1(\Omega)\) by (F1); along a subsequence they
converge weakly in \(H^1\) and strongly in \(L^2\) to functions
\(v_1,\ldots,v_j\). The extension mass inside the shrinking hole tends
to zero by the uniform \(L^4\) bound and its vanishing area. Therefore
the limits are orthonormal.

On a fixed outer collar the trace converges strongly in
\(L^2(\partial\Omega)\), by compactness of the trace map. The
inner-boundary quadratic term tends to zero by (F3), even when
\(\alpha<0\). For any fixed coefficient vector \(z\), weak lower
semicontinuity on \(\Omega\setminus\overline B_\rho\), followed by
\(\rho\downarrow0\), gives

\[
 q_0\left(\sum_{l=1}^jz_lv_l,\sum_{l=1}^jz_lv_l\right)
 \le\liminf q_\varepsilon\left(\sum_{l=1}^jz_lv_l^\varepsilon,
                                    \sum_{l=1}^jz_lv_l^\varepsilon\right).
 \tag{F6}
\]

Here the mass and outer-boundary terms converge; excluding a fixed disk
before taking the limit avoids charging the artificial
extension\textquotesingle s gradient energy to the original form. Choose
the subsequence to realize the lower limit of the \(j\)th eigenvalue.
The right side of (F6) is at most
\(\liminf(\lambda_j(\Omega_\varepsilon)+c)\,|z|^2\). Minmax on the
\(j\)-dimensional limiting span proves the lower bound. This argument
also proves convergence along the full sequence. \(\square\)

For any isolated limiting eigenvalue, its full finite-dimensional
cluster is consequently separated by a uniform gap from the
complementary spectrum. This is the only spectral convergence input
needed below. Interior analyticity follows from the constant-coefficient
equation \(-\Delta u=\lambda u\), regardless of the regularity of the
outer boundary.

\Needspace{8\baselineskip}\section{5. The normalized exterior
problem}\label{the-normalized-exterior-problem}

Consider the Hilbert space

\[
 X=\{v\in H^1(D\cap B_S)\text{ for every sufficiently large }S:
       \nabla v\in L^2(D)\}/\mathbb R,
 \qquad \|[v]\|_X=\|\nabla v\|_{L^2(D)}.
\]

Completeness follows by selecting representatives with zero average on
\(A=B_R\setminus\overline\Sigma\). Poincaré inequalities on each
connected \(B_S\setminus\overline\Sigma\), with this fixed average, turn
a gradient-Cauchy sequence into an \(H^1\)-Cauchy sequence on every
bounded truncation.

If \(g\in L^2(\partial\Sigma)\) has integral zero, the functional

\[
 [v]\longmapsto\int_{\partial\Sigma}gv
\]

is well defined and bounded on \(X\): subtract the average on \(A\), and
apply the trace and Poincaré inequalities on \(A\). Lax-Milgram gives a
unique class \([W]\) satisfying

\[
 \int_D\nabla W\cdot\nabla v
 =\int_{\partial\Sigma}gv\qquad([v]\in X).
 \tag{S7}
\]

For clarity, compactly supported \(H^1\) test functions are dense in
this gradient norm. First truncate a representative at heights
\(\pm M\); the gradient error tends to zero as \(M\to\infty\). For fixed
\(M\), multiply by a radial cutoff equal to one on \(B_L\), zero outside
\(B_{L^2}\), and linear in \(\log r\) between these radii. Its squared
gradient integral is \(2\pi/\log L\). The cutoff error tends to zero by
boundedness of the truncated function and the integrability of its
gradient. Local smoothing in Lipschitz charts gives the corresponding
density for smooth tests up to the boundary.

The solution is harmonic outside a fixed disk. Its Fourier expansion
there has the constant mode, a possible logarithmic mode, and the modes
\(r^{\pm j}\cos(j\theta)\) and \(r^{\pm j}\sin(j\theta)\). Finite
Dirichlet energy excludes the logarithmic and growing modes. The
remaining expansion shows that a unique representative satisfies

\[
 W(y)=O(|y|^{-1}),\qquad \nabla W(y)=O(|y|^{-2}).
 \tag{S8}
\]

One may obtain the estimates by expanding on a circle outside the hole
and using convergence of the decaying Fourier series beyond twice that
radius. This removes the additive constant, which cannot be controlled
by a homogeneous Sobolev inequality in dimension two.

For \(g=\partial_\nu P_k\), the compatibility condition follows from

\[
 \int_{\partial\Sigma}\partial_\nu P_k
 =-\int_\Sigma\Delta P_k=0.
\]

Thus (S1) exists and is unique. Its energy is
\(E_k=\int_{\partial\Sigma}(\partial_\nu P_k)W_k\). It is nonzero:
otherwise \(W_k\) is constant, hence \(g=0\); testing \(g\) against the
trace of \(P_k\) would give

\[
 0=\int_{\partial\Sigma}(\partial_\nu P_k)P_k
   =-\int_\Sigma|\nabla P_k|^2,
\]

impossible for a nonconstant polynomial on a nonempty open set. In
particular \(C_k>0\).

\Needspace{8\baselineskip}\section{6. Quantitative cutoff
correctors}\label{quantitative-cutoff-correctors}

Choose \(c>0\) so that, for all small \(\varepsilon\),

\[
 q_\varepsilon(v,w)=
 \int_{\Omega_\varepsilon}(\nabla v\cdot\nabla w+cvw)
 +\alpha\int_{\partial\Omega_\varepsilon}vw
\]

is uniformly equivalent to the \(H^1\) inner product. This follows from
the fixed outer-boundary trace inequality and the small-hole trace
estimate in Section 4 (see also {[}FRS26{]}). Set \(\mu=\lambda+c>0\).
For every fixed \(\eta\in(0,1/2)\), the latter estimate is

\[
 \|v\|_{L^2(\partial(\varepsilon\Sigma))}
 \le C_\eta\varepsilon^{1/2-\eta}\|v\|_{H^1(\Omega_\varepsilon)}.
 \tag{S9}
\]

Define \(V_\varepsilon\) by

\[
 q_\varepsilon(V_\varepsilon,v)=L_\varepsilon(v),\qquad
 L_\varepsilon(v)=
 \int_{\partial(\varepsilon\Sigma)}(\partial_\nu u+\alpha u)v.
 \tag{S10}
\]

The explicit approximate corrector is

\[
 A_\varepsilon(x)=\varepsilon^k\chi(x)W_k(x/\varepsilon).
 \tag{S11}
\]

The cutoff equals one on the entire hole for small \(\varepsilon\). On
its boundary, the discrepancy between the data in (S10) and
\(\partial_\nu A_\varepsilon+\alpha A_\varepsilon\) has \(L^2\) norm
\(O(\varepsilon^{k+1/2})\). Indeed, the leading normal derivative is
exactly \(\varepsilon^{k-1}\partial_\nu P_k\), the Taylor remainder is
\(O(\varepsilon^k)\), \(u=O(\varepsilon^k)\) there, and the trace of
\(W_k\) on \(\partial\Sigma\) belongs to \(L^2\).

From (S8), scaling and the compact support of \(\chi\),

\[
 \|A_\varepsilon\|_{L^2(\Omega_\varepsilon)}
 \le C\varepsilon^{k+1}\sqrt{1+|\log\varepsilon|},
 \qquad
 \|\nabla A_\varepsilon\|_{L^2(\Omega_\varepsilon)}=O(\varepsilon^k).
 \tag{S12}
\]

On the support of \(\nabla\chi\), both
\(\varepsilon^kW_k(x/\varepsilon)\) and its \(x\)-gradient are
\(O(\varepsilon^{k+1})\). Since \(W_k\) is harmonic,

\[
 \|(-\Delta+c)A_\varepsilon\|_{L^2(\Omega_\varepsilon)}
 \le C\varepsilon^{k+1}\sqrt{1+|\log\varepsilon|}.
\]

Integration by parts, (S9), and the boundary discrepancy therefore give

\[
 |L_\varepsilon(v)-q_\varepsilon(A_\varepsilon,v)|
 \le C_\eta\varepsilon^{k+1-\eta}
       \|v\|_{H^1(\Omega_\varepsilon)}.
\]

The boundary integration is valid weakly for the exterior solution; no
\(H^2\) regularity at a Lipschitz corner is required. Uniform coercivity
and (S10) imply

\[
 \|V_\varepsilon-A_\varepsilon\|_{H^1(\Omega_\varepsilon)}
 \le C_\eta\varepsilon^{k+1-\eta}.
 \tag{S13}
\]

In particular

\[
 \|V_\varepsilon\|_{L^2}=O_\eta(\varepsilon^{k+1-\eta}),\qquad
 T_\varepsilon:=q_\varepsilon(V_\varepsilon,V_\varepsilon)
 =\varepsilon^{2k}E_k+O(\varepsilon^{2k+1}).
 \tag{S14}
\]

For the energy estimate, the discarded far-field gradient energy in
(S11) is \(O(\varepsilon^{2k+2})\); the mass term is
\(O(\varepsilon^{2k+2}(1+|\log\varepsilon|))\); the inner-boundary term
is \(O(\varepsilon^{2k+1})\); and the outer-boundary term vanishes. For
the cross term, write \(R_\varepsilon=V_\varepsilon-A_\varepsilon\) and
use the exact identity

\[
 q_\varepsilon(R_\varepsilon,A_\varepsilon)
 =L_\varepsilon(A_\varepsilon)-q_\varepsilon(A_\varepsilon,A_\varepsilon).
\]

The boundary residual paired with \(A_\varepsilon\) is
\(O(\varepsilon^{2k+1})\) by their boundary \(L^2\) norms. The bulk
residual paired with \(A_\varepsilon\) is
\(O(\varepsilon^{2k+2}(1+|\log\varepsilon|))\) by (S12). Finally
\(q_\varepsilon(R_\varepsilon,R_\varepsilon)
=O_\eta(\varepsilon^{2k+2-2\eta})=o(\varepsilon^{2k+1})\). This
establishes the sharper error in (S14) without discarding cancellation
in the cross term.

\Needspace{8\baselineskip}\section{7. Simple-eigenvalue reduction and
normalization}\label{simple-eigenvalue-reduction-and-normalization}

Write \(f_\varepsilon=u-V_\varepsilon\), with \(u\) restricted to the
perforated domain. The weak equation and (S10) give

\[
 q_\varepsilon(f_\varepsilon,v)=\mu\langle u,v\rangle,
 \qquad
 (\mathcal A_\varepsilon-\mu)f_\varepsilon=\mu V_\varepsilon,
 \tag{S15}
\]

where \(\mathcal A_\varepsilon\) is the positive operator associated
with \(q_\varepsilon\). The second equality holds in \(L^2\), since the
first places \(f_\varepsilon\) in its operator domain.

Spectral convergence and simplicity give a uniform gap between \(\mu\)
and every eigenvalue of \(\mathcal A_\varepsilon\) except
\(\mu_\varepsilon=\lambda_\varepsilon+c\). Let \(\Pi_\varepsilon\) be
the orthogonal projection onto that eigenspace, and put
\(r_\varepsilon=(I-\Pi_\varepsilon)f_\varepsilon\). Expanding (S15) in
an orthonormal eigenbasis gives

\[
 \|r_\varepsilon\|_{H^1}
 \le C\|V_\varepsilon\|_{L^2}.
 \tag{S16}
\]

The bound follows because
\(\sup_{j\ne n}\mu_j^\varepsilon/(\mu_j^\varepsilon-\mu)^2\) is
uniformly finite, including the large eigenvalues.

The useful exact identity is

\[
 \mu\langle u,V_\varepsilon\rangle
 =L_\varepsilon(u)-T_\varepsilon,
 \tag{S17}
\]

obtained by interchanging the two arguments in \(q_\varepsilon\) and
using (S10). Further,

\[
 L_\varepsilon(u)=
 -\int_{\varepsilon\Sigma}(|\nabla u|^2-\lambda u^2)
 +\alpha\int_{\partial(\varepsilon\Sigma)}u^2.
 \tag{S18}
\]

Hence \(L_\varepsilon(u)=O(\varepsilon^{2k})\) and
\(\langle u,V_\varepsilon\rangle=O(\varepsilon^{2k})\). Consequently

\[
 \|f_\varepsilon\|_{L^2}^2
 =1-\int_{\varepsilon\Sigma}u^2
    -2\langle u,V_\varepsilon\rangle+\|V_\varepsilon\|_{L^2}^2
 =1+O(\varepsilon^{2k}).
\]

Equations (S14) and (S16) show
\(\|\Pi_\varepsilon f_\varepsilon\|_{L^2}=1+O(\varepsilon^{2k})\).
Normalize this projection with its positive overlap with \(u\). Its
\(H^1\) norm is uniformly bounded, so

\[
 \|u_\varepsilon-f_\varepsilon\|_{H^1}
 =O_\eta(\varepsilon^{k+1-\eta})+O(\varepsilon^{2k})
 =O_\eta(\varepsilon^{k+1-\eta}).
 \tag{S19}
\]

This proves (S4) by (S13). It also avoids an unjustified passage from a
global little-oh estimate to local rescaled \(L^2\) convergence.

For the eigenvalue, calculate using (S15) and (S17):

\[
 q_\varepsilon(f_\varepsilon,f_\varepsilon)
   -\mu\|f_\varepsilon\|_{L^2}^2
 = L_\varepsilon(u)-T_\varepsilon-\mu\|V_\varepsilon\|_{L^2}^2.
\]

The contribution of \(r_\varepsilon\) to the left side has absolute
value \(O(\|V_\varepsilon\|_{L^2}^2)\) by (S16). Thus

\[
 (\lambda_\varepsilon-\lambda)
       \|\Pi_\varepsilon f_\varepsilon\|_{L^2}^2
 =L_\varepsilon(u)-T_\varepsilon
       +O(\|V_\varepsilon\|_{L^2}^2).
 \tag{S20}
\]

Taylor expansion in (S18) yields

\[
 L_\varepsilon(u)
 =-\varepsilon^{2k}\int_\Sigma|\nabla P_k|^2
    +O(\varepsilon^{2k+1}).
\]

Insert this and (S14) in (S20). The extra errors are
\(O_\eta(\varepsilon^{2k+2-2\eta})\) and \(O(\varepsilon^{4k})\); both
are smaller than the error in (S3). This proves (S3).

For (S5), uniform extension into the hole and planar Sobolev embedding
imply, for fixed \(S\),

\[
 \|v(\varepsilon\,\cdot)\|_{L^2(B_S\setminus\overline\Sigma)}
 \le C_{\eta,S}\varepsilon^{-\eta}
             \|v\|_{H^1(\Omega_\varepsilon)}.
 \tag{S21}
\]

Indeed, apply Hölder with exponent \(p=2/\eta\) on a set of area
\(O(\varepsilon^2)\) after extension. Applying (S21) to the error in
(S4) and dividing by \(\varepsilon^k\) gives
\(O(\varepsilon^{1-2\eta})\). The rescaled gradient error is
\(O(\varepsilon^{1-\eta})\). The Taylor expansion completes (S5).
Finally, (S4), (S12), and the energy calculation following (S14) prove
(S6).

\Needspace{8\baselineskip}\section{8. The second nodal
coefficient}\label{the-second-nodal-coefficient}

Let \(P_{k+1}\) denote the homogeneous Taylor term of degree \(k+1\) in
\(u\), allowing \(P_{k+1}=0\). It is harmonic, since the right side of
\(-\Delta u=\lambda u\) begins at degree \(k\). Denote its decaying
exterior corrector by \(W_{k+1}\), with the same definition as (S1). Put

\[
 \mathcal C(P_k,P_{k+1})=
 \int_\Sigma\nabla P_k\cdot\nabla P_{k+1}
 +\int_D\nabla W_k\cdot\nabla W_{k+1}.
\]

The two-term refinement is

\[
 \begin{aligned}
 \lambda_\varepsilon-\lambda
 ={}&-C_k\varepsilon^{2k}\\
 &+\varepsilon^{2k+1}\left[
    \alpha\int_{\partial\Sigma}(P_k-W_k)^2
       -2\mathcal C(P_k,P_{k+1})\right]
    +o(\varepsilon^{2k+1}).
 \end{aligned}
 \tag{S22}
\]

Here is the remainder calculation. With \(A_\varepsilon\) and
\(R_\varepsilon\) as above, (S13) gives
\(q_\varepsilon(R_\varepsilon,R_\varepsilon)=o(\varepsilon^{2k+1})\).
Consequently the exact identity

\[
 T_\varepsilon
 =2L_\varepsilon(A_\varepsilon)
       -q_\varepsilon(A_\varepsilon,A_\varepsilon)
       +q_\varepsilon(R_\varepsilon,R_\varepsilon)
\]

can be expanded to the next order using only the known corrector
\(W_k\):

\[
 \begin{aligned}
 L_\varepsilon(A_\varepsilon)
 ={}&\varepsilon^{2k}E_k
 +\varepsilon^{2k+1}
   \int_{\partial\Sigma}(\partial_\nu P_{k+1}+\alpha P_k)W_k
 +O(\varepsilon^{2k+2}),\\
 q_\varepsilon(A_\varepsilon,A_\varepsilon)
 ={}&\varepsilon^{2k}E_k
 +\alpha\varepsilon^{2k+1}\int_{\partial\Sigma}W_k^2
 +O(\varepsilon^{2k+2}(1+|\log\varepsilon|)).
 \end{aligned}
\]

Meanwhile (S18) gives

\[
 \begin{aligned}
 L_\varepsilon(u)
 ={}&-\varepsilon^{2k}\int_\Sigma|\nabla P_k|^2\\
 &+\varepsilon^{2k+1}
   \left[-2\int_\Sigma\nabla P_k\cdot\nabla P_{k+1}
               +\alpha\int_{\partial\Sigma}P_k^2\right]
   +O(\varepsilon^{2k+2}).
 \end{aligned}
\]

The \(L^2\) spectral error in (S20) is
\(O_\eta(\varepsilon^{2k+2-2\eta})=o(\varepsilon^{2k+1})\); the
normalization error is \(O(\varepsilon^{4k})=o(\varepsilon^{2k+1})\).
Finally, (S7) implies

\[
 \int_{\partial\Sigma}(\partial_\nu P_{k+1})W_k
 =\int_D\nabla W_{k+1}\cdot\nabla W_k.
\]

Substituting these identities into (S20) gives (S22).

If \(\Sigma=-\Sigma\), uniqueness and the far-field normalization imply
\(W_j(-y)=(-1)^jW_j(y)\). The mixed term in (S22) then vanishes by
parity. For a centrally symmetric hole the next coefficient is therefore
exactly \(\alpha\int_{\partial\Sigma}(P_k-W_k)^2\).

\Needspace{8\baselineskip}\section{9. The nonnodal logarithmic
coefficient}\label{the-nonnodal-logarithmic-coefficient}

Now assume \(u_0:=u(0)\ne0\) and \(\alpha\ne0\). Let

\[
 p=\mathcal H^1(\partial\Sigma),\qquad
 h(y)=\nabla u(0)\cdot\nu(y)+\alpha u_0,
 \qquad m=\int_{\partial\Sigma}h=\alpha u_0p.
\]

Choose a point \(y_*\in\Sigma\) and put \(F(y)=-(m/(2\pi))\log|y-y_*|\).
The normal convention gives \(\int_{\partial\Sigma}\partial_\nu F=m\).
The zero-mean argument of Section 5 supplies a decaying harmonic \(W\)
with boundary data \(h-\partial_\nu F\). Thus \(U=F+W\) solves
\(\partial_\nu U=h\) and

\[
 U(y)=-\frac m{2\pi}\log|y|+O(|y|^{-1}).
\]

Take the cutoff approximation

\[
 B_\varepsilon(x)=\varepsilon\chi(x)
 \left[U(x/\varepsilon)-\frac m{2\pi}\log\varepsilon\right].
 \tag{S23}
\]

The subtraction is essential: on a fixed physical annulus the bracket is
\(-(m/(2\pi))\log|x|+O(\varepsilon/|x|)\), with no growing constant. It
follows by direct integration that

\[
 \|B_\varepsilon\|_{L^2}=O(\varepsilon),\qquad
 q_\varepsilon(B_\varepsilon,B_\varepsilon)
 =\frac{m^2}{2\pi}\varepsilon^2\log(1/\varepsilon)
      +O(\varepsilon^2).
 \tag{S24}
\]

To check the remainder in the energy: the harmonic logarithm contributes
\(m^2/(2\pi)\) per unit logarithmic radial length; its cross integral
with the decaying part has a convergent tail; the fixed cutoff region
contributes \(O(\varepsilon^2)\); and the inner Robin term is
\(O(\varepsilon^3\log^2\varepsilon)=o(\varepsilon^2)\).

The bulk residual of (S23) is \(O(\varepsilon)\) in \(L^2\). Its
boundary discrepancy in (S10) has \(L^2\) norm
\(O(\varepsilon^{3/2}(1+|\log\varepsilon|))\). Estimate (S9) makes its
dual norm \(o(\varepsilon)\), for any fixed \(\eta<1/2\). Coercivity
gives

\[
 \|V_\varepsilon-B_\varepsilon\|_{H^1}=O(\varepsilon),\qquad
 \|V_\varepsilon\|_{L^2}=O(\varepsilon).
\]

For the energy cross term, write
\(S_\varepsilon=V_\varepsilon-B_\varepsilon\) and use
\(q_\varepsilon(S_\varepsilon,B_\varepsilon)
=L_\varepsilon(B_\varepsilon)-q_\varepsilon(B_\varepsilon,B_\varepsilon)\).
The bulk residual pairs with \(B_\varepsilon\) to give
\(O(\varepsilon^2)\); the boundary pairing is
\(O(\varepsilon^3(1+|\log\varepsilon|)^2)\). Also
\(q_\varepsilon(S_\varepsilon,S_\varepsilon)=O(\varepsilon^2)\).
Therefore

\[
 T_\varepsilon
 =\frac{m^2}{2\pi}\varepsilon^2\log(1/\varepsilon)
   +O(\varepsilon^2).
 \tag{S25}
\]

Equations (S16)-(S20) remain valid. This time
\(L_\varepsilon(u)=\alpha p u_0^2\varepsilon+O(\varepsilon^2)\),
\(\langle u,V_\varepsilon\rangle=O(\varepsilon)\), and
\(\|\Pi_\varepsilon f_\varepsilon\|_{L^2}^2=1+O(\varepsilon)\). It
follows that

\[
 \lambda_\varepsilon-\lambda
 =\alpha p u_0^2\varepsilon
 -\frac{\alpha^2p^2u_0^2}{2\pi}\varepsilon^2\log(1/\varepsilon)
 +O(\varepsilon^2).
 \tag{S26}
\]

The first term is established in {[}FRS26{]}. Earlier small-Robin-hole
and logarithmic-expansion results must also be compared before a
priority claim for the second term.

\Needspace{8\baselineskip}\section{10. Disk and ellipse
coefficients}\label{disk-and-ellipse-coefficients}

For \(\Sigma=B_r\) and
\(P_k(\rho,\theta)=\rho^k(A\cos k\theta+B\sin k\theta)\),

\[
 W_k(\rho,\theta)=-r^{2k}\rho^{-k}
                         (A\cos k\theta+B\sin k\theta),
 \qquad
 C_k=2k\pi r^{2k}(A^2+B^2).
 \tag{S27}
\]

The coefficients use the homogeneous Taylor polynomial itself, including
its factorial normalization.

For an ellipse with semiaxes \(a,b>0\), use exterior conformal
coordinates

\[
 z=\frac{a+b}{2}\left(w+\frac{a-b}{a+b}w^{-1}\right),\qquad |w|>1.
\]

For \(P_1=x\), the exterior solution is \(W_1=-b\rho^{-1}\cos\theta\) in
\(w=\rho e^{i\theta}\) coordinates. For \(P_1=y\) it is
\(W_1=-a\rho^{-1}\sin\theta\). The common conformal metric factor
cancels from the normal-data equality; Dirichlet energy is conformally
invariant. Thus for \(P_1=\xi x+\zeta y\),

\[
 C_1=\pi b(a+b)\xi^2+\pi a(a+b)\zeta^2.
 \tag{S28}
\]

This is a useful check that arbitrary-hole coefficients need not depend
only on area. It is not counted as an independent discovery or paper.

\Needspace{8\baselineskip}\section{11. Multiple eigenvalues and their vanishing
orders}\label{multiple-eigenvalues-and-their-vanishing-orders}

Use the geometry and normal convention of Section 3 and the estimates of
Section 4. Fix \(\alpha\ne0\). Let \(\lambda\) be an eigenvalue of
multiplicity \(m\) on \(\Omega\), and let \(E\) be its real eigenspace
with the \(L^2(\Omega)\) inner product. Write
\(p=\mathcal H^1(\partial\Sigma)\) and

\[
 E^+=\{u\in E:u(0)=0\}.
\]

If evaluation at zero is nonzero on \(E\), choose the unit vector
\(u_0\) orthogonal to \(E^+\) with \(a=u_0(0)>0\). Otherwise set
\(E^+=E\) and omit \(u_0\).

For integers \(k\ge1\), let \(F_k\subset E^+\) consist of functions
whose derivatives of every order below \(k\) vanish at zero. Thus
\(F_1=E^+\). Set \(H_k=F_k\ominus F_{k+1}\), using the \(L^2\)
orthogonal complement inside \(F_k\). Discard zero spaces and enumerate
the remaining degrees as

\[
 1\le k_1<\cdots<k_s,\qquad d_r=\dim H_{k_r}.
\]

This list is finite and \(E^+=\bigoplus_{r=1}^s H_{k_r}\). Indeed, the
descending finite-dimensional filtration eventually stabilizes; its
intersection is zero by interior analyticity and unique continuation.
Any nonzero vector of \(H_{k_r}\) has exactly order \(k_r\), and the map
taking it to its homogeneous leading Taylor polynomial is injective. In
two dimensions this gives \(d_r\le2\).

Choose an orthonormal basis \(u_i\) adapted to these spaces. Denote its
orders by \(\kappa_i\) (in increasing groups) and its leading harmonic
polynomials by \(P_i\). For a nonconstant harmonic polynomial \(P\), let
\(W(P)\) be the decaying exterior solution of

\[
 \Delta W=0\ \text{in }D,\qquad
 \partial_\nu W=\partial_\nu P\ \text{on }\partial\Sigma.
\]

Define the bilinear form and its matrix

\[
 \mathcal C(P,Q)=\int_\Sigma\nabla P\cdot\nabla Q
                 +\int_D\nabla W(P)\cdot\nabla W(Q),
 \qquad C_{ij}=\mathcal C(P_i,P_j).
 \tag{M1}
\]

The matrix \(C\) is positive definite. To see this, a zero quadratic
form forces \(\sum_i t_iP_i\) to be constant on an open set, hence
constant everywhere. Its homogeneous terms have positive degrees, so it
is zero. Within each degree, injectivity of the leading-polynomial map
gives all \(t_i=0\).

Partition \(C\) by the groups \(H_{k_r}\). Let \(C_{<r,<r}\) denote its
principal block containing all smaller orders. Define

\[
 S_1=C_{11},\qquad
 S_r=C_{rr}-C_{r,<r}C_{<r,<r}^{-1}C_{<r,r}\quad(r>1).
 \tag{M2}
\]

Every \(S_r\) is positive definite. Its eigenvalues
\(s_{r,1},\ldots,s_{r,d_r}\) are independent of the orthonormal bases
chosen in each \(H_{k_r}\).

\Needspace{6\baselineskip}\textbf{Theorem 11.1 (all nodal branches).} The perturbed cluster admits
a labeling with one branch for each pair \((r,j)\), satisfying

\[
 \lambda_{r,j}(\varepsilon)
 =\lambda-s_{r,j}\varepsilon^{2k_r}
       +o(\varepsilon^{2k_r}).
 \tag{M3}
\]

If evaluation on \(E\) is identically zero these are all \(m\) branches.
Otherwise they are \(m-1\) branches, and the remaining branch satisfies

\[
 \lambda_0(\varepsilon)=\lambda+\alpha p a^2\varepsilon+o(\varepsilon).
 \tag{M4}
\]

Thus for sufficiently small holes, every nodal branch decreases. If
evaluation is nonzero, positive \(\alpha\) gives exactly one increasing
branch and \(m-1\) decreasing branches; negative \(\alpha\) gives \(m\)
decreasing branches. These assertions concern shifts relative to the
limiting \(\lambda\), not monotonicity as a function of hole size. For
\(E^+=0\), only (M4) occurs. No claim of complete simplicity is made
when a block \(S_r\) has repeated eigenvalues.

The all-nodal assertion also holds for \(\alpha=0\), using Neumann
spectral convergence. The mixed Neumann case is excluded: its nonnodal
branch has a different scale and can couple with degree-one data at
leading order.

\Needspace{8\baselineskip}\section{12. Weighted reduction for a full
cluster}\label{weighted-reduction-for-a-full-cluster}

Fix a coercive shift \(c\) and let \(\mu=\lambda+c>0\). All scalar
products below are over \(\Omega_\varepsilon\). For any \(u\in E\) put

\[
 L_u(v)=\int_{\partial(\varepsilon\Sigma)}
             (\partial_\nu u+\alpha u)v,
 \qquad q_\varepsilon(V_u,v)=L_u(v),\qquad f_u=u-V_u.
 \tag{M5}
\]

Let \(\Pi_\varepsilon\) be the orthogonal projection onto the full
\(m\)-dimensional cluster converging to \(\lambda\), for the shifted
operator. Define \(g_u=\Pi_\varepsilon f_u\) and \(r_u=f_u-g_u\).
Spectral convergence gives a uniform gap to the complementary spectrum.
As in Section 7,

\[
 (\mathcal A_\varepsilon-\mu)f_u=\mu V_u,
 \qquad \|r_u\|_{H^1}\le C\|V_u\|_{L^2}.
 \tag{M6}
\]

Since \(f_u\to u\) and \(r_u\to0\) in \(L^2\) uniformly on the unit
sphere of the fixed finite-dimensional space \(E\), \(u\mapsto g_u\) is
an isomorphism onto the cluster for small \(\varepsilon\).

In the selected basis, the cluster shifts are exactly the generalized
eigenvalues of the symmetric pair

\[
 A_{ij}=q_\varepsilon(g_{u_i},g_{u_j})-\mu\langle g_{u_i},g_{u_j}\rangle,
 \qquad G_{ij}=\langle g_{u_i},g_{u_j}\rangle,
 \qquad G\longrightarrow I.
 \tag{M7}
\]

Write \(b_\varepsilon=q_\varepsilon-\mu\langle\cdot,\cdot\rangle\). The
following exact identity, rather than an unweighted norm estimate, keeps
the information at every vanishing order:

\[
 A_{ij}=L_{u_i}(u_j)-q_\varepsilon(V_{u_i},V_{u_j})
       -\mu\langle V_{u_i},V_{u_j}\rangle
       -b_\varepsilon(r_{u_i},r_{u_j}).
 \tag{M8}
\]

To verify it, symmetry and the eigenfunction equation give

\[
 \mu\langle u_j,V_{u_i}\rangle
 =L_{u_i}(u_j)-q_\varepsilon(V_{u_i},V_{u_j}).
\]

Expand \(b_\varepsilon(f_{u_i},f_{u_j})\) and subtract the complementary
spectral contribution. The cross terms between the cluster and its
complement vanish for both \(q_\varepsilon\) and the mass inner product.
Also

\[
 L_{u_i}(u_j)
 =-\int_{\varepsilon\Sigma}
       (\nabla u_i\cdot\nabla u_j-\lambda u_i u_j)
      +\alpha\int_{\partial(\varepsilon\Sigma)}u_i u_j.
 \tag{M9}
\]

For a nodal basis vector define

\[
 Z_i=\varepsilon^{\kappa_i}\chi W(P_i)(\cdot/\varepsilon),
 \qquad R_i=V_{u_i}-Z_i.
\]

Section 6, whose corrector estimates did not use simplicity, gives for
fixed \(0<\eta<1/2\)

\[
 \|Z_i\|_{H^1}=O(\varepsilon^{\kappa_i}),\quad
 \|R_i\|_{H^1}=O_\eta(\varepsilon^{\kappa_i+1-\eta}),\quad
 \|V_{u_i}\|_{L^2}=O_\eta(\varepsilon^{\kappa_i+1-\eta}).
 \tag{M10}
\]

Scaling the two leading terms in (M8), and using Cauchy-Schwarz on their
remainders, yields for any two nodal indices

\[
 A_{ij}=-\varepsilon^{\kappa_i+\kappa_j}
             (C_{ij}+o(1)).
 \tag{M11}
\]

For example the energy cross remainder is
\(O_\eta(\varepsilon^{\kappa_i+\kappa_j+1-\eta})\); the mass and
complementary terms in (M8) are
\(O_\eta(\varepsilon^{\kappa_i+\kappa_j+2-2\eta})\). The finitely many
entrywise errors in (M11) are uniform. In particular, with
\(D_\varepsilon=\operatorname{diag}(\varepsilon^{\kappa_i})\), the whole
nodal block is

\[
 A_{++}=-D_\varepsilon(C+o(1))D_\varepsilon.
 \tag{M12}
\]

This stronger statement cannot be replaced by a single unweighted
\(o(\varepsilon^{2k_1})\) bound.

\Needspace{8\baselineskip}\section{13. The nonnodal direction in a multiple
cluster}\label{the-nonnodal-direction-in-a-multiple-cluster}

Assume \(u_0\) is present. The logarithmic corrector in Section 9 gives

\[
 \|V_{u_0}\|_{L^2}=O(\varepsilon),\qquad
 \|V_{u_0}\|_{H^1}=O(\varepsilon\sqrt{1+|\log\varepsilon|}),
 \qquad A_{00}=\alpha p a^2\varepsilon+o(\varepsilon).
 \tag{M13}
\]

For a nodal index \(j\), equation (M9) gives
\(L_{u_0}(u_j)=O(\varepsilon^{\kappa_j+1})\). For the mixed corrector
energy, use its defining equation to obtain the sharper bound

\[
 q_\varepsilon(V_{u_0},V_{u_j})
 =L_{u_0}(Z_j)+q_\varepsilon(V_{u_0},R_j)
 =O(\varepsilon^{\kappa_j+1}).
 \tag{M14}
\]

Indeed the first boundary integral has data of order one, trace
\(\varepsilon^{\kappa_j}W(P_j)\) and length of order \(\varepsilon\).
The second term is
\(O_\eta(\varepsilon^{\kappa_j+2-\eta}\sqrt{1+|\log\varepsilon|})\),
which is smaller. The remaining two terms in (M8) are bounded by
\(O_\eta(\varepsilon^{\kappa_j+2-\eta})\). Hence

\[
 A_{0j}=O(\varepsilon^{\kappa_j+1}).
 \tag{M15}
\]

Put \(d=A_{00}\) and \(v=(A_{0j})\). For small \(\varepsilon\),
\(d\ne0\). The change of basis

\[
 T_0=\begin{pmatrix}1&-d^{-1}v\\0&I\end{pmatrix}
\]

tends to the identity and gives the exact congruence

\[
 T_0^TAT_0=
 \begin{pmatrix}d&0\\0&A_{++}-d^{-1}v^Tv\end{pmatrix}.
 \tag{M16}
\]

The correction in its nodal block has entries
\(O(\varepsilon^{\kappa_i+\kappa_j+1})\). Thus it does not change the
limiting matrix \(C\) in (M12), for either sign of \(d\). The
transformed mass matrix still tends to \(I\).

\Needspace{8\baselineskip}\section{14. Schur complements and signed
shifts}\label{schur-complements-and-signed-shifts}

For completeness we prove the finite-dimensional step, in a block form
of the method already supplied by {[}ALM26{]}. Let
\(C(\varepsilon)\to C>0\). Gaussian block elimination in increasing
order gives an upper block-triangular matrix \(U(\varepsilon)\), with
identity diagonal blocks, for which

\[
 U(\varepsilon)^TC(\varepsilon)U(\varepsilon)
  =\operatorname{diag}(S_1(\varepsilon),\ldots,S_s(\varepsilon)),
 \qquad S_r(\varepsilon)\longrightarrow S_r.
 \tag{M17}
\]

All matrices depend continuously on \(C(\varepsilon)\) because its
principal blocks are uniformly positive definite. Completing the square
in the preceding blocks shows that the limiting \(S_r\) is exactly (M2).

Set \(T_1=D_\varepsilon^{-1}U(\varepsilon)D_\varepsilon\). Its block
above the diagonal, in row \(r\) and column \(t>r\), is
\(\varepsilon^{k_t-k_r}U_{rt}(\varepsilon)\), so \(T_1\to I\). It
transforms the nodal matrix in (M12), or in (M16), exactly into

\[
 -\operatorname{diag}\big(
   \varepsilon^{2k_1}S_1(\varepsilon),\ldots,
   \varepsilon^{2k_s}S_s(\varepsilon)\big).
 \tag{M18}
\]

After both transformations the mass matrix is \(I+o(1)\), whereas the
stiffness matrix is the direct sum of (M18) and, if present, \(d\). The
ordinary eigenvalues of this block matrix have exactly the asymptotics
(M3)-(M4).

Here is a justification that replacing \(I\) by \(I+o(1)\) preserves
those asymptotics even when the matrix is indefinite. For a positive
definite mass matrix \(B\) with \((1-\delta)I\le B\le(1+\delta)I\),
consider a negative eigenvalue \(h_j<0\) of a real symmetric matrix
\(H\). The number of generalized eigenvalues below \(-t\) is the
negative inertia of \(H+tB\). Since

\[
 H+t(1-\delta)I\le H+tB\le H+t(1+\delta)I,
\]

monotonicity of eigenvalues in the matrix order, or its minmax proof,
bounds the corresponding negative generalized eigenvalue \(\theta_j\) by

\[
 \frac{h_j}{1-\delta}\le\theta_j\le\frac{h_j}{1+\delta}.
 \tag{M19}
\]

Multiplicity is handled by allowing equality and then approaching
thresholds from either side. Applying the same statement to \(-H\)
bounds the positive eigenvalues. Congruence preserves inertia, so there
is no unaccounted change in their numbers. Letting \(\delta\to0\) proves
multiplicative equivalence separately for every nonzero eigenvalue,
however small its scale. This proves Theorem 11.1.

\Needspace{8\baselineskip}\section{15. Common-order clusters at second
order}\label{common-order-clusters-at-second-order}

Suppose every nonzero \(u\in E\) has the same positive order \(k\). Thus
evaluation vanishes identically, \(s=1\), and \(m\le2\). Write

\[
 u_i=P_i+Q_i+O(|x|^{k+2}),
\]

where \(P_i\) and \(Q_i\) are the homogeneous Taylor polynomials of
degree \(k\) and \(k+1\). Both are harmonic, and \(Q_i\) may be zero.
Define

\[
 B_{ij}=\alpha\int_{\partial\Sigma}
        (P_i-W(P_i))(P_j-W(P_j))
       -\mathcal C(P_i,Q_j)-\mathcal C(Q_i,P_j).
 \tag{M20}
\]

\Needspace{6\baselineskip}\textbf{Theorem 15.1 (common-order second coefficient).} For each
eigenvalue \(c_*\) of \(C\), let \(\mathcal E_{c_*}\) be its Euclidean
eigenspace, and let \(b_1,\ldots,b_q\) be the eigenvalues of the
compression of \(B\) to \(\mathcal E_{c_*}\). The corresponding branches
satisfy

\[
 \lambda_j(\varepsilon)
 =\lambda-c_*\varepsilon^{2k}
           +b_j\varepsilon^{2k+1}+o(\varepsilon^{2k+1}).
 \tag{M21}
\]

To prove this, use the bilinear energy identity

\[
 q_\varepsilon(V_{u_i},V_{u_j})
 =L_{u_i}(Z_j)+L_{u_j}(Z_i)
       -q_\varepsilon(Z_i,Z_j)+q_\varepsilon(R_i,R_j).
 \tag{M22}
\]

Its last term is
\(O_\eta(\varepsilon^{2k+2-2\eta})=o(\varepsilon^{2k+1})\); expanding
the first three terms gives, exactly as in the two-term calculation of
Section 8,

\[
 A=\varepsilon^{2k}(-C+\varepsilon B+o(\varepsilon)).
 \tag{M23}
\]

In more detail, the terms containing \(Q_i,Q_j\) are the two mixed
interior energies from (M9) plus their exterior energies, using
\(\int_{\partial\Sigma}(\partial_\nu Q_i)W(P_j)
 =\int_D\nabla W(Q_i)\cdot\nabla W(P_j)\). The Robin terms combine as
the boundary square in (M20). Bulk mass, cutoff and complementary
projection terms are all smaller than \(\varepsilon^{2k+1}\).

We also need more than \(G\to I\) here. The exact identity preceding
(M9) gives \(\langle u_i,V_{u_j}\rangle=O(\varepsilon^{2k})\). Expanding
\(\langle f_{u_i},f_{u_j}\rangle\) and subtracting
\(\langle r_{u_i},r_{u_j}\rangle\) gives

\[
 G=I+O(\varepsilon^{2k})=I+o(\varepsilon).
 \tag{M24}
\]

Conjugating (M23) by \(G^{-1/2}\) therefore leaves its first two
coefficients unchanged. Finally diagonalize \(C\). In a fixed eigenspace
of \(C\), the first perturbation is the compression of \(B\); the
coupling to the orthogonal complement contributes \(O(\varepsilon^2)\)
after a block elimination across the fixed spectral gaps of \(C\).
Equivalently, the Schur complement of \(-C+\varepsilon B-zI\) at
\(z=-c_*+\varepsilon b\) is
\(\varepsilon(B|_{\mathcal E_{c_*}}-bI)+o(\varepsilon)\). Continuity of
the eigenvalues of these finite symmetric matrices gives (M21).

For a centrally symmetric hole, the two mixed terms in (M20) vanish by
parity. This does not automatically split a repeated \(c_*\): the
compressed boundary form must itself have distinct eigenvalues.

\Needspace{8\baselineskip}\section{16. Analytic dependence on hole
shape}\label{analytic-dependence-on-hole-shape}

For a harmonic polynomial \(P\) without a constant term, let \(W_{t,P}\)
solve the exterior problem of Section 5 on
\(D_t=\mathbb R^2\setminus\overline\Sigma_t\). Thus \(W_{t,P}\) decays
at infinity and its normal derivative into the hole equals that of
\(P\). Define

\[
 \mathcal C_t(P,Q)=\int_{\Sigma_t}\nabla P\cdot\nabla Q
       +\int_{D_t}\nabla W_{t,P}\cdot\nabla W_{t,Q}.
 \tag{H2}
\]

For each fixed pair \(P,Q\), this is a real-analytic function of \(t\)
on a sufficiently small fixed interval \((-t_*,t_*)\).

Here are details that address the two-dimensional constant ambiguity.
Choose a smooth diffeomorphism \(\Phi_t\) of the plane, mapping \(D_0\)
to \(D_t\), equal to the radial boundary displacement in a collar of the
circle, and equal to the identity outside \(B_R\). The area-normalized
displacement is real analytic in \(t\) with coefficients smooth in
space; a radial cutoff supplies such a map, uniformly invertible after
decreasing \(t_*\). Pull back the weak problem to

\[
 X=\{v\in H^1_{\mathrm{loc}}(\overline D_0):\nabla v\in L^2(D_0)\}/\mathbb R.
\]

The notation means \(H^1\) on every bounded exterior truncation,
including its inner boundary. Its norm is the gradient norm. The
pulled-back form has coefficient

\[
 A_t=|\det D\Phi_t|\,(D\Phi_t)^{-1}(D\Phi_t)^{-T}.
\]

It is uniformly coercive and real analytic as a bounded operator on
\(X\). The boundary functional is

\[
 \ell_{t,P}(v)=\int_{\partial\Sigma_t}
                    (\partial_{\nu_t}P)(v\circ\Phi_t^{-1})\,ds.
\]

It has integral-zero data, since \(P\) is harmonic. Subtracting a fixed
annular average and using the trace bound makes it a bounded functional
on \(X\). In the pulled-back boundary parametrization its coefficients
are analytic in \(t\) in every finite smooth norm, so it is analytic in
\(X^*\). The inverse of the coercive operator is locally analytic by its
norm-convergent Neumann series. This proves analytic dependence of the
solution class.

Because \(\Phi_t\) is the identity outside \(B_R\), every representative
is harmonic there. Subtract its mean on one fixed exterior circle.
Finite Dirichlet energy excludes growing modes and a logarithm; this
subtraction selects exactly the decaying representative. The circle mean
is a bounded functional on anchored representatives, so this
normalization preserves analytic dependence on bounded truncations. The
energy in (H2) can be written as the pulled-back bilinear form applied
to the solution classes; the interior integral follows by change of
variables. This proves the asserted analyticity of \(\mathcal C_t\)
without a planar homogeneous Sobolev inequality.

We will differentiate a boundary equation at \(t=0\). This is justified
by pulling it back to the fixed smooth circle: the differentiated weak
solution satisfies a smooth, uniformly elliptic problem in a fixed
collar. Difference quotients converge in \(H^1\), and local Neumann
regularity upgrades this convergence in each fixed finite Sobolev order
after applying the differentiated equation. Equivalently one may
differentiate the coefficients and boundary data in the pulled-back
variational equation first. Outside \(B_R\) the Fourier coefficients are
continuous linear functionals of the solution traces and therefore
analytic as well.

\Needspace{8\baselineskip}\section{17. Shape variation at the
disk}\label{shape-variation-at-the-disk}

It suffices to differentiate the unnormalized radial family: its
normalization changes the shape only at order \(t^2\). Put

\[
 U_{t,P}=P-W_{t,P}.
\]

This has zero normal derivative on the hole, with
\(U_{t,P}-P=O(|x|^{-1})\) at infinity. For

\[
 P(r,\theta)=r^k p(\theta),\qquad
 p(\theta)=a\cos k\theta+b\sin k\theta,
\]

the disk solution is

\[
 U_{0,P}(r,\theta)=(r^k+r^{-k})p(\theta).
 \tag{H3}
\]

Let \(\dot U_P\) be the Eulerian derivative at zero. It is harmonic for
\(r>1\) and decays at infinity. Differentiate the zero normal derivative
on the radial graph. Multiplying by the harmless normalizing factor,
that equation reads

\[
 \partial_r U_{t,P}(1+tf,\theta)
   -\frac{t f'}{(1+tf)^2}\partial_\theta U_{t,P}(1+tf,\theta)=0.
\]

Using \(\partial_{rr}U_{0,P}(1,\theta)=2k^2p(\theta)\) and
\(\partial_\theta U_{0,P}(1,\theta)=2p'(\theta)\) gives

\[
 \partial_r\dot U_P(1,\theta)
       =-2k^2 f p+2f'p'=2(fp')'.
 \tag{H4}
\]

The right-hand side has mean zero, so it has a unique decaying exterior
Neumann solution. Notice that the radial derivative here points out of
the hole; the zero condition itself is independent of that sign choice.

We now relate a decaying multipole to (H2), to avoid assuming a
shape-derivative formula. If \(Q\) is harmonic and \(P=r^k p\),
Green\textquotesingle s identity on \(B_R\setminus\overline\Sigma_t\)
yields

\[
 \mathcal C_t(P,Q)
 =\int_{\partial B_R}
        (P\partial_r W_{t,Q}-W_{t,Q}\partial_rP)\,ds.
 \tag{H5}
\]

Indeed the corresponding inner-boundary integral is
\(-\int_{\Sigma_t}\nabla P\cdot\nabla Q
 -\int_{D_t}\nabla W_{t,P}\cdot\nabla W_{t,Q}\). The second equality
uses the exterior variational identity; all functions have the required
zero-mean data and finite energy. If the \(k\)th decaying Fourier mode
of \(W_{t,Q}\) is \(r^{-k}w_k(\theta)\), (H5) becomes

\[
 \mathcal C_t(P,Q)=-2k\int_0^{2\pi}p(\theta)w_k(\theta)\,d\theta.
 \tag{H6}
\]

For \(P\) and \(Q\) of the same degree \(k\), write \(Q=r^kq(\theta)\).
The \(k\)th Fourier mode of \(\dot U_Q\) on the circle is
\(-k^{-1}\operatorname{proj}_k[2(fq')']\), by (H4). Since
\(\dot W_Q=-\dot U_Q\), differentiation of (H6) and integration by parts
give

\[
 \left.\frac{d}{dt}\mathcal C_t(P,Q)\right|_{t=0}
   =-4\int_0^{2\pi}p(fq')'
   =4\int_0^{2\pi}f p'q'.
 \tag{H7}
\]

This is the familiar tangential-energy shape variation in this setting.
It is derived here for the needed disk case; shape differentiation of
generalized polarization tensors is established prior literature and is
not claimed new.

Only the constant and \(2k\) Fourier components of \(f\) contribute to
(H7). For a zero-mean \(f\) with \(2k\) component
\(A_k\cos2k\theta+B_k\sin2k\theta\), the derivative matrix in the
polynomial basis \((r^k\cos k\theta,r^k\sin k\theta)\) is

\[
 2\pi k^2
 \begin{pmatrix}-A_k&-B_k\\-B_k&A_k\end{pmatrix}.
 \tag{H8}
\]

For (H1), \(A_k=e^{-k^2}>0\) and \(B_k=0\). Its two eigenvalues are
opposite and nonzero, for every positive integer \(k\).

\Needspace{8\baselineskip}\section{18. Complete splitting of every Schur
block}\label{complete-splitting-of-every-schur-block}

Fix an unperforated eigenspace \(E\) and use the orthogonal
vanishing-order decomposition of Section 11. This basis is fixed in
\(\Omega\) and does not depend on \(t\). Let \(C(t)\) be the matrix with
entries \(\mathcal C_t(P_i,P_j)\), and let \(S_r(t)\) be its Schur
blocks, eliminating smaller positive orders first. Analyticity and
positive definiteness of \(C(t)\) imply that every \(S_r(t)\) is
positive definite and analytic on the whole small interval under
consideration.

At \(t=0\) the matrix \(C(0)\) has no cross terms between distinct
degrees, by orthogonality of Fourier modes on the disk and on its
exterior. Consequently

\[
 S_r(0)=C_{rr}(0),\qquad S_r'(0)=C_{rr}'(0).
 \tag{H9}
\]

The second identity follows by differentiating
\(C_{rr}-C_{r,<r}C_{<r,<r}^{-1}C_{<r,r}\): both outer factors of the
correction vanish at zero.

A block of dimension one is already simple. For a block of dimension
two, let \(k\) be its order and write its two leading polynomials in the
standard harmonic basis using an invertible real coefficient matrix
\(M\). At the disk,

\[
 C_{rr}(0)=2k\pi M^TM.
 \tag{H10}
\]

If this matrix has distinct eigenvalues, they stay distinct for
sufficiently small \(t\). Otherwise it equals \(cI\) for some \(c>0\),
so \(M=\sqrt{c/(2k\pi)}\,O\) with \(O\) orthogonal. Equations (H8)-(H10)
then give

\[
 S_r'(0)=c k e^{-k^2}\,
           O^T\begin{pmatrix}-1&0\\0&1\end{pmatrix}O.
 \tag{H11}
\]

The eigenvalue difference of \(S_r(t)\) is therefore

\[
 2c k e^{-k^2}|t|+o(|t|).
 \tag{H12}
\]

In particular, it is positive for every sufficiently small nonzero
\(t\). There are finitely many blocks in \(E\), so all become simple for
every sufficiently small nonzero \(t\). The nonnodal direction, when
present, is one-dimensional and moves at order \(\varepsilon\), while
all positive-order blocks move at distinct even powers of
\(\varepsilon\). Thus it cannot collide with them for sufficiently small
\(\varepsilon\). Within a block, distinct positive Schur coefficients
give distinct leading eigenvalue shifts by Theorem 11.1. This proves the
assertion for a fixed eigenspace and every sufficiently small nonzero
\(t\).

For a two-dimensional block define its discriminant

\[
 \Delta_r(t)=(\operatorname{tr}S_r(t))^2-4\det S_r(t).
 \tag{H13}
\]

It is analytic and not identically zero, by (H10) or (H12). Its zero set
in \((-t_*,t_*)\) is discrete, hence countable. Taking the finite union
over the blocks of an eigenspace and then the countable union over the
discrete unperforated spectrum gives a countable exceptional set
\(\mathcal N\). For every \(t\) outside that set, every Schur block of
every eigenspace has simple coefficients. Apply Theorem 11.1 separately
to each cluster. For a finite spectral window, take the minimum of the
finitely many thresholds and also ensure that different limiting
clusters remain separated by their positive limiting gaps. This proves
Theorem 2.1.

If a finite window ends inside a multiple limiting cluster, include that
entire cluster when taking this finite minimum. The result then implies
simplicity for the requested first \(M\) eigenvalues.

\Needspace{8\baselineskip}\section{19. An exact conformal coefficient
calculation}\label{an-exact-conformal-coefficient-calculation}

There is a useful independent coefficient check for every positive
integer \(k\). Consider the exterior conformal map

\[
 z=\Psi_t(w)=w+t w^{-(2k-1)},\qquad |w|>1.
 \tag{H14}
\]

For sufficiently small real \(t\), the boundary is a smooth simple
perturbation of the circle and the map is univalent on its exterior. For
example, the boundary is embedded and the derivative has no zeros for
small \(t\); the argument principle, together with the simple pole at
infinity, gives degree one on the unbounded complementary component. Its
first normal displacement is \(\cos(2k\theta)\), though the boundary
parametrization also has a tangential displacement. Its enclosed area is
\(\pi(1-(2k-1)t^2)\) by the boundary area integral.

For \(P=\operatorname{Re}z^k\) one has

\[
 \Psi_t(w)^k=w^k+k t w^{-k}
           +\sum_{j=2}^k\binom{k}{j}t^j w^{-(2j-1)k}.
 \tag{H15}
\]

In these conformal coordinates the zero-Neumann total field is
\(\operatorname{Re}(w^k+w^{-k})\), since every growing mode is \(w^k\)
and each boundary radial derivative must vanish. Therefore the
pulled-back corrector is

\[
 W=\operatorname{Re}\left((kt-1)w^{-k}
           +\sum_{j=2}^k\binom{k}{j}t^j w^{-(2j-1)k}\right).
\]

The exterior Dirichlet energy is the sum of \(\pi n\) times the squared
decaying Fourier coefficients. Green\textquotesingle s identity on the
hole gives the interior energy as \(\pi k\) minus that same weighted sum
for the negative coefficients in (H15), namely the sum with \(kt\)
rather than \(kt-1\) at the first mode. Adding the two energies cancels
every higher mode and gives the exact formula

\[
 \mathcal C_t(\operatorname{Re}z^k,\operatorname{Re}z^k)
       =2\pi k(1-kt).
 \tag{H16}
\]

For \(P=\operatorname{Im}z^k\), the negative Laurent powers have the
opposite sine sign on the boundary. The same calculation gives

\[
 \mathcal C_t(\operatorname{Im}z^k,\operatorname{Im}z^k)
       =2\pi k(1+kt),\qquad
 \mathcal C_t(\operatorname{Re}z^k,\operatorname{Im}z^k)=0.
 \tag{H17}
\]

Scaling the hole to area \(\pi\) multiplies both coefficients by
\((1-(2k-1)t^2)^{-k}\). Thus the two derivatives at zero are exactly
\(-2\pi k^2\) and \(2\pi k^2\), agreeing with (H8). This is an exact
harmonic calculation, not a numerical PDE test. This one-mode conformal
family is used to check a fixed degree; the all-degree construction in
(H1) remains the family used in the theorem.

\Needspace{8\baselineskip}\section{20. Scope, source comparison and review
status}\label{scope-source-comparison-and-review-status}

The main assertion is one coherent result about a fixed shrinking-hole
family. Its simple-eigenvalue formulas, multiple-eigenvalue reduction,
shape derivative and splitting consequence are not counted as separate
papers. The mixed nonnodal Neumann case is excluded from Theorem 11.1.
Complete simplicity at one common positive hole size for the whole
infinite spectrum is not proved. No parameter-uniform or
geometry-uniform threshold is asserted.

The source comparison uses the final June 2026 Robin article
{[}FRS26{]}, the inspected arXiv versions of {[}FLO25, FLO26, ALM26{]},
and the openly accessible older Robin-hole statements {[}Oza92, OR92{]}.
The 2017 report {[}Dab17{]}, including its variational principle and all
application statements, was also inspected. It is cited separately from
{[}Dab22{]}, whose final text was unavailable. The full text of
{[}AKLZ10, AKLZ12, LdC12, Dab22, Gad93{]} has not been completely
compared in this investigation. Their abstracts or primary indexed
excerpts establish relevant prior frameworks but do not suffice to
exclude theorem-level overlap. No absence from a search is treated as
proof of novelty.

The originating assistant has checked the internal derivations,
including the signs in the reduced matrix, the weighted remainders, the
logarithmic normalization, and the countable-exception argument. This is
not independent mathematical review. Markdown formula checks establish
syntax only. The accompanying source audit records precisely which
sources were read and which retrievals failed. The contribution is
presented as one unreviewed preprint. The mathematical argument is a
hand proof; the finite diagnostics are supplementary and are not proof
dependencies. No guarantee of exhaustive priority or journal acceptance
is made.

\raggedright\Needspace{8\baselineskip}\section{References}\label{references}

{[}FRS26{]} V. Felli, P. Roychowdhury and G. Siclari, \emph{Quantitative
spectral stability for the Robin Laplacian}, Mathematische Annalen 395,
article 96 (2026).
\href{https://doi.org/10.1007/s00208-026-03523-4}{Final article};
\href{https://arxiv.org/abs/2504.05994v1}{arXiv:2504.05994v1}.

{[}FLO25{]} V. Felli, L. Liverani and R. Ognibene, \emph{Quantitative
spectral stability for the Neumann Laplacian in domains with small
holes}, Journal of Functional Analysis 288 (2025), 110817.
\href{https://doi.org/10.1016/j.jfa.2024.110817}{DOI};
\href{https://arxiv.org/abs/2311.05999v3}{arXiv:2311.05999v3}.

{[}FLO26{]} V. Felli, L. Liverani and R. Ognibene, \emph{On the
splitting of Neumann eigenvalues in perforated domains},
\href{https://arxiv.org/abs/2601.13129v1}{arXiv:2601.13129v1}, 2026.

{[}ALM22{]} L. Abatangelo, C. Léna and P. Musolino, \emph{Ramification
of multiple eigenvalues for the Dirichlet-Laplacian in perforated
domains}, Journal of Functional Analysis 283 (2022), 109718.
\href{https://doi.org/10.1016/j.jfa.2022.109718}{DOI}. See the
coefficient correction {[}ALM26{]}.

{[}ALM26{]} L. Abatangelo, C. Léna and P. Musolino, \emph{Corrigendum to
``Ramification of multiple eigenvalues for the Dirichlet-Laplacian in
perforated domains''},
\href{https://arxiv.org/abs/2606.03996v1}{arXiv:2606.03996v1}, 2026.

{[}AKLZ10{]} H. Ammari, H. Kang, M. Lim and H. Zribi, \emph{Layer
potential techniques in spectral analysis. Part I: Complete asymptotic
expansions for eigenvalues of the Laplacian in domains with small
inclusions}, Transactions of the American Mathematical Society 362
(2010), 2901-2922.
\href{https://doi.org/10.1090/S0002-9947-10-04695-7}{DOI}.

{[}AKLZ12{]} H. Ammari, H. Kang, M. Lim and H. Zribi, \emph{The
generalized polarization tensors for resolved imaging. Part I: Shape
reconstruction of a conductivity inclusion}, Mathematics of Computation
81 (2012), 367-386.
\href{https://people.math.inha.ac.kr/~hbkang/paper/AKLZ_PT_july20.pdf}{Author-hosted
copy}.

{[}LdC12{]} M. Lanza de Cristoforis, \emph{Simple Neumann eigenvalues
for the Laplace operator in a domain with a small hole. A functional
analytic approach}, Revista Matemática Complutense 25 (2012), 369-412.
\href{https://doi.org/10.1007/s13163-011-0081-8}{DOI}.

{[}Dab17{]} A. Dabrowski, \emph{A variational principle for the
perturbation of repeated eigenvalues and applications}, ETH Seminar for
Applied Mathematics Research Report 2017-37, July 2017.
\href{https://www.sam.math.ethz.ch/sam_reports/reports_final/reports2017/2017-37.pdf}{Repository
PDF}.

{[}Dab22{]} A. D\k{a}browski, \emph{Asymptotic behavior of repeated
eigenvalues of perturbed self-adjoint elliptic operators}, Asymptotic
Analysis 126 (2022), 187-202.
\href{https://doi.org/10.3233/ASY-201669}{DOI}.

{[}Gad93{]} R. R. Gadyl\textquotesingle shin, \emph{Splitting a multiple
eigenvalue in the boundary value problem for a membrane clamped on a
small part of the boundary}, Siberian Mathematical Journal 34 (1993),
433-450.
\href{https://www.mathnet.ru/php/archive.phtml?jrnid=smj&option_lang=eng&paperid=1651&wshow=paper}{Primary
record}.

{[}Oza92{]} S. Ozawa, \emph{Singular variation of domain and spectra of
the Laplacian with small Robin conditional boundary I}, Osaka Journal of
Mathematics 29 (1992), 837-850.
\href{https://ir.library.osaka-u.ac.jp/repo/ouka/all/7188/ojm29_04_10.pdf}{Repository
PDF}.

{[}OR92{]} S. Ozawa and S. Roppongi, \emph{Singular variation of domain
and spectra of the Laplacian with small Robin conditional boundary II},
Kodai Mathematical Journal 15 (1992), 403-429.
\href{https://www.jstage.jst.go.jp/article/kodaimath1978/15/3/15_3_403/_pdf}{Journal
PDF}.

\end{document}
